% sec-01-introduction.tex - Section 1, Introduction.  DRAFT stage.
% One page maximum of prose.  Two figure slots, Figures 1 and 2, and two tables,
% Tables 4 and 5.  Every drafting instruction names the exact bib file or
% repository directory stage 7 must read.
\section{Introduction}
\label{sec:introduction}

Between January 2025 and July 2026 the executive branch committed to artificial
intelligence and to cancer at the same time, and in the same documents. Executive
Order 14179 removed barriers to American leadership in AI \cite{eo14179}, and
Executive Order 14212 established the MAHA Commission \cite{eo14212}. Executive
Order 14355, signed September 30, 2025, is the only executive order aimed
squarely at both topics: it tasks the MAHA Commission, in coordination with HHS,
the Assistant to the President for Science and Technology, and the Special
Advisor for AI and Crypto, with developing ways to use AI for improved
diagnoses, treatments, cures, and prevention of pediatric cancer
\cite{eo14355, eo14355app, akin_eo14355}.

\draftinstr{Stage 7: expand this section to one full page using
\appfile{new-trial-system/references/trump-ai-cancer-2025-2026.bib}. The
required citation keys, in order, are eo14179, eo14212, eo14355 with
eo14355app and akin\_eo14355, ccdidoubling\_nih, eo14363 with
doe\_genesis\_launch, biogenesis\_nih with hpcwire\_genesis,
crs\_ai\_healthcare, hhs\_ai\_strategy, cms\_healthtech, maha\_strategy,
trialblazer, aacr\_policymonitor\_2026, cancerdiscovery\_fy27 and
science\_fundingbill. Tone: the administration and the FDA are ambitious about
cancer and about large language models, and this paper's system is the sponsor
side answer to that ambition. Do not exceed one page of prose.}

\figslot{mermaid-type / flowchart LR}{7.6}{Figure 1. Eleven Federal AI and cancer
actions from January 2025 to July\\2026, the three capabilities they supply, and
the one they leave unfilled.}

The same period funded the commitment. HHS raised the Childhood Cancer Data
Initiative budget from \$50 million to \$100 million and committed to bringing
private-sector partners in to apply advanced AI \cite{ccdidoubling_nih}.
Executive Order 14363 launched the Genesis Mission \cite{eo14363}, and its NIH
arm, the Bio Genesis Mission announced July 22, 2026, carries more than \$1.2
billion in obligated FY2026 and planned FY2027 funding
\cite{biogenesis_nih, hpcwire_genesis}. Table 4 places these actions on one
chronology.

\draftinstr{Stage 7: build Table 4 from the same bib file. Eleven to fourteen
rows, columns Action, Date, What it supplies, Citation key. Include the
June 22, 2026 HHS department-wide clinical trial effort and Operation
TrialBlazer.}

\begin{apptable}
\begin{tabularx}{\textwidth}{>{\raggedright\arraybackslash}p{5.4cm} >{\raggedright\arraybackslash}p{2.1cm} >{\raggedright\arraybackslash}X}
\headrow \thead{Action} & \thead{Date} & \thead{What it supplies to a pancreatic cancer sponsor}\\
\midrule
EO 14179, Removing Barriers to American Leadership in AI & Jan 23, 2025 & Policy cover for AI in a regulated submission \\
EO 14212, MAHA Commission & Feb 13, 2025 & The parent authority EO 14355 cites \\
FDA Elsa generative AI tool & Jun 2, 2025 & An agency that reads AI-assisted protocol text \\
EO 14355, Unlocking Cures for Pediatric Cancer with AI & Sep 30, 2025 & The only order aimed at AI and cancer together \\
CCDI budget doubled to \$100 million & Sep 30, 2025 & Money for AI-driven oncology data work \\
EO 14363, Launching the Genesis Mission & Nov 24, 2025 & An integrated Federal AI platform \\
HHS AI Strategy and the agentic AI platform & Dec 1 to 4, 2025 & Agency-side capacity to process AI-assisted filings \\
Operation TrialBlazer & Jun 22, 2026 & Faster trial activation at cancer centers \\
HHS department-wide clinical trial effort & Jun 22, 2026 & A stated intent to restore U.S. trial leadership \\
Bio Genesis Mission & Jul 22, 2026 & Over \$1.2 billion toward AI in biology \\
\end{tabularx}
\end{apptable}
\vspace{-0.6cm}
\tabcap{Table 4. Ten Federal actions between January 2025 and July 2026, and what\\each
supplies to a sponsor preparing a pancreatic cancer trial dossier.}

What none of them supplies is a sponsor-side method. Compute, data, platform,
agency review capacity and faster site activation are all being built; the
document a sponsor must file is still written the way it was written in 2015.
That gap is the subject of this paper, and Figure 1 states it directly.

\draftinstr{Stage 7: expand the gap argument, then introduce the paper's four
deliverables in one sentence each, keyed to sections 3 through 7.}

\figslot{d2-type / grid}{7.2}{Figure 2. Ten operating axes on which the prior and
new trial systems differ,\\each cell carrying a measured value from the author's
deposited record.}

Table 5 and Figure 2 carry the same ten axes: the first as text a reader can
cite, the second as a grid a reader can scan. The prior system is not described
here as bad work. It is described as work whose unit is a team-month, and a
team-month cannot produce an IND, two protocols, five bill versions and fourteen
funding artifacts inside a single quarter.

\draftinstr{Stage 7: build Table 5 from
\appfile{new-trial-system/d2/fig-02-old-versus-new-system-grid.md}, cell value
table. Ten rows, columns Axis, Prior system, New system, Source. The figure and
the table must agree exactly, value for value.}

\begin{apptable}
\begin{tabularx}{\textwidth}{>{\raggedright\arraybackslash}p{3.5cm} >{\raggedright\arraybackslash}p{3.9cm} >{\raggedright\arraybackslash}p{3.9cm} >{\raggedright\arraybackslash}X}
\headrow \thead{Operating axis} & \thead{Prior system} & \thead{New system} & \thead{Source}\\
\midrule
Document production unit & Team-month & Prompt-hour & Build record \\
IND assembly time & Months across a group & Four days, one author & \S3 \\
Peer review entry point & After completion & During development & \S7 \\
Review latency & 7 to 8 weeks best case & Same day & \S7 \\
Reviewers per round & 2 to 3 humans & 3 model manufacturers & \S7 \\
Virtual trial cost & Above \$120,000 per run & \$36,330 projected & \S6 \\
Provenance record & Acknowledgement paragraph & Commit history plus DOI & \S2 \\
Figure format & Raster, not reusable & Machine readable, reusable & \S2 \\
Revision granularity & Whole manuscript & One file, one commit & \S2 \\
Author count & Group, plus a CRO & One, plus the models & Every deposit \\
\end{tabularx}
\end{apptable}
\vspace{-0.6cm}
\tabcap{Table 5. The ten operating axes of Figure 2 as citable text, with the\\section
of this paper where each value is established.}

\draftinstr{Stage 8: check that this section still fits one page of prose after
the two floats settle, and that Figure 1, Figure 2, Table 4 and Table 5 are each
referenced once by name above.}
